\subsection*{Texto E. Convocatoria}
\textit{Texto original de práctica. Institución y actividad ficticias.}

La Biblioteca del Parque invita a residentes mayores de dieciséis años al taller gratuito de crónica vecinal. Habrá veinte lugares. Las solicitudes se recibirán del 3 al 10 de mayo mediante el formulario disponible en recepción. Cada solicitante entregará un texto propio de entre 300 y 500 palabras sobre un lugar de su colonia. No se requiere experiencia previa.

Si se reciben más solicitudes válidas que lugares, se realizará un sorteo el 12 de mayo. La lista de personas seleccionadas se publicará el día 13. Las sesiones se celebrarán los sábados 16, 23 y 30 de mayo de 10 a 12 horas. Quienes no puedan asistir a las tres sesiones deberán comunicarlo antes del sorteo para liberar su lugar.

\subsection*{Texto F. Nota informativa}
\textit{Texto original de práctica. Hechos y lugar ficticios.}

El lunes, la Dirección de Servicios de Villa Clara anunció una prueba de separación de residuos en cuatro calles del centro. Durante seis semanas, un vehículo recogerá materiales reciclables los miércoles y otro mantendrá la recolección habitual los demás días. La prueba no modifica el servicio en otras colonias.

Según el comunicado, el municipio medirá el peso de los materiales separados y registrará las incidencias del recorrido. La directora indicó que todavía no se ha decidido ampliar el programa. Al concluir la prueba, se publicará un informe y se abrirá un periodo de comentarios vecinales. La asociación del mercado pidió que ese informe incluya también el costo de operación; el comunicado inicial no ofreció esa cifra.
